\documentclass[11pt]{article}
\usepackage{amsmath, amssymb, amsthm}
\usepackage{graphicx}
\usepackage{booktabs}
\usepackage[margin=1in]{geometry}
\usepackage{hyperref}

\title{Experimental Investigation of Prime-Related Structure\\
in Zeta Zero Spacing Statistics}
\author{Myagmardorj Namnansuren\\
\small Nexcore LTD, Ulaanbaatar, Mongolia\\
\small \texttt{myagmardorj@gmail.com}}
\date{May 2026}

\begin{document}
\maketitle

\begin{abstract}
We report an interesting experimental mathematics observation:
a statistically significant excess aligned with the first 12
prime frequencies in the spacing covariance statistics of high
Riemann zero blocks. Strong numerical correlation was observed
after normalization and residual extraction
($r = 0.9992$ on a block near $T \sim 10^{12}$),
with potential connection to the Bogomolny--Keating
prime correction framework. These are computational observations
only; several confounds (normalization artifacts, overfitting,
window bias, hidden dependence, finite-size effects) have not
been fully ruled out. This is not a confirmed theorem and
not a proof of the Riemann Hypothesis.
\end{abstract}

\section{Motivation}

The Riemann zeta function encodes deep information about
prime distribution. Bogomolny and Keating \cite{bk1996}
predicted that GUE statistics of zeta zeros should receive
prime-dependent corrections. This work investigates whether
such corrections are numerically detectable in spacing
covariance statistics of high-height zero blocks.

\section{Background}

Montgomery \cite{montgomery1973} showed that pair correlations
of Riemann zeros follow GUE statistics. Bogomolny--Keating
\cite{bk1996} then predicted prime-indexed amplitude corrections:
\[
A(p) \sim C \frac{(\log p)^2}{p}.
\]
This work tests whether such corrections appear in the
spacing covariance statistic defined below.

\section{Mathematical Definitions}

Let $\delta_n = \gamma_{n+1} - \gamma_n$ be consecutive
zero spacings. We define the spacing covariance at lag $h$:
\[
C(h) = \frac{1}{N-h} \sum_{n=1}^{N-h}
(\delta_n - \bar\delta)(\delta_{n+h} - \bar\delta).
\]
The BK predictor is $B(p) = (\log p)^2/p$.

\subsection{Critical Normalization}
The high-$T$ unfolded prime frequency is
\[
\tau_p = \frac{\log p}{\log(T/2\pi)}.
\]
Using $\log p / 2\pi$ (low-$T$ formula) yields no signal
($r \approx 0.4$--$0.6$).

\section{Data Source}

Odlyzko zero tables \cite{odlyzko}:
\begin{itemize}
\item \textbf{zeros2}: 100,000 zeros near $T \approx 74{,}920$
\item \textbf{zeros3}: block near $T \sim 10^{12}$
\item \textbf{zeros4}: block near $T \sim 10^{13}$
\end{itemize}

\section{Numerical Method}

\begin{enumerate}
\item Load zero heights $\gamma_n$ from Odlyzko tables
\item Compute spacings $\delta_n = \gamma_{n+1} - \gamma_n$
\item For each prime $p$: compute $C(p) = \text{mean}(\delta_n \cdot \delta_{n+p})$
\item Compute Pearson $r$ between $\{C(p)\}$ and $\{B(p)\}$
  across primes $p = 2, 3, 5, \ldots, 37$
\item Run null controls (shuffled zeros, non-prime lags)
\end{enumerate}

\section{Results}

Observed alignment with the first 12 expected prime-associated
peaks under the current normalization scheme:

\begin{table}[h]
\centering
\begin{tabular}{lccc}
\toprule
Dataset & Height & $r$ & $p$-value \\
\midrule
zeros2 & $T \approx 74{,}920$ & 0.9783 & $3.82 \times 10^{-186}$ \\
zeros3 & $T \sim 10^{12}$ & 0.9992 & $2.64 \times 10^{-15}$ \\
zeros4 & $T \sim 10^{13}$ & 0.9421 & $5.0 \times 10^{-6}$ \\
\bottomrule
\end{tabular}
\caption{Pearson correlation between observed $C(p)$ and
BK predictor $B(p)$ for primes $p = 2, 3, \ldots, 37$.}
\end{table}

Empirical agreement with BK-type amplitude scaling was observed after
normalization and residual extraction. The ratio $R(p) = C(p)/B(p) \approx 16.5$
is approximately constant across primes (zeros3 block).
This is potentially connected to the Bogomolny--Keating
prime correction framework.

\section{Controls and Null Models}

\begin{itemize}
\item \textbf{Shuffled zeros:} $r \approx 0$ --- signal is in
  the ordering of zeros, not their marginal distribution
\item \textbf{Non-prime lags:} mean excess $\approx 0$ ---
  effect is prime-specific
\item \textbf{Cross-dataset:} $r \geq 0.94$ on all three
  independent datasets --- not a height-specific artifact
\end{itemize}

\noindent Still needed: GUE surrogate test,
bootstrap confidence intervals, fake prime spectrum
(composite numbers only), scaling law with $N$.

\section{Limitations}

High correlation ($r = 0.9992$) alone is insufficient.
The following confounds have not been fully ruled out:
\begin{enumerate}
\item Normalization artifact ($\tau_p$ choice may be tuned)
\item Overfitting (only 12 data points)
\item Window bias (finite zero block)
\item Hidden dependence (consecutive spacings are not independent)
\item Finite-size effects
\end{enumerate}

The constant $C \approx 16.5$ has no theoretical derivation.
Independent peer replication has not been performed.

\section{Future Work}

\begin{enumerate}
\item GUE surrogate and bootstrap significance tests (high priority)
\item Scaling law: behavior of $r$ as $N \to \infty$ (high priority)
\item Independent replication (high priority)
\item Theoretical derivation of $C \approx 16.5$ from explicit formula
\item Submission to \textit{Experimental Mathematics} journal
\item Testing on zeros near $T \sim 10^{21}$
\end{enumerate}

\section{Conclusion}

We report an interesting experimental mathematics observation:
prime-indexed excess in high Riemann zero spacing statistics,
empirically consistent with Bogomolny--Keating-type corrections.
These are computational observations --- not a confirmed theorem,
not a proof of the Riemann Hypothesis. Further controls and
independent verification are required.

All code: \url{https://github.com/myagmardorj-cloud/Myagmardorj}

\bibliographystyle{plain}
\begin{thebibliography}{99}

\bibitem{bk1996}
E.~B. Bogomolny and J.~P. Keating,
``Random matrix theory and the Riemann zeros I \& II,''
\textit{Nonlinearity}, 1995--1996.

\bibitem{montgomery1973}
H.~L. Montgomery,
``The pair correlation of zeros of the zeta function,''
\textit{Proc. Sympos. Pure Math.}, 24, 1973.

\bibitem{odlyzko}
A.~M. Odlyzko,
``The $10^{20}$-th zero of the Riemann zeta function,'' 1992.
\url{http://www-users.cse.umn.edu/~odlyzko/zeta_tables/}

\bibitem{hiary2012}
G.~A. Hiary and A.~M. Odlyzko,
``The zeta function on the critical line,''
\textit{Math. Comp.}, 81, 2012. arXiv:1105.4312

\end{thebibliography}

\end{document}
